% =============================================================================
% 1. INTRODUCTION
% =============================================================================
\section{Introduction}
\label{sec:introduction}

The global maintenance, repair and overhaul (\MRO{}) market is valued at over
USD~90 billion annually and remains a critical determinant of airline
operating costs, fleet availability and on-time performance
\citep{Zhang2025,Boyd2015}. The sector operates under one of the most
stringent safety regimes across all engineering domains, where
maintenance-attributable errors continue to account for between 12\% and 18\%
of commercial aviation accidents \citep{Boyd2015,FAA2024} despite decades of
regulatory and technological progress. The economic and environmental
consequences of these errors extend well beyond the immediate corrective
action: a single warranty claim arising from a defective overhaul has been
reported to exceed USD~189{,}500 in direct material and labour costs alone
\citep{Pereira2017}, and the indirect cost --- expressed as material scrap,
energy consumed in rework, displaced operational capacity, missed
slot-allocation opportunities at congested hubs and reputational damage ---
has been estimated at three to seven times the direct figure
\citep{Zhang2025}. For airlines, every maintenance disruption translates into
unit cost pressure and competitive position erosion; for the \MRO{} operator,
each event compounds into quality-cost dilution and the slow attrition of
preferred-supplier status.

Safety and sustainability in \MRO{} are therefore intrinsically linked.
Maintenance errors do not only generate accident risk; they also generate
waste. Every component scrapped because of a procedural deviation, every
engine returned for rework, every additional litre of jet fuel consumed by a
delayed turnaround and every kilowatt-hour of energy invested in corrective
actions correspond to a sustainability loss that is rarely reported in
isolation but is systematically captured by quality-cost accounting
\citep{Chen2023,Yang2024}. This connection has become particularly relevant
under the rising pressure of stakeholder demand for environmental, social and
governance (\ESG{}) disclosures from airlines and their supply chain
\citep{Sousa2025}.

Within this context, the technical training of maintenance personnel emerges
as a high-leverage intervention. Training is mandated by both the civil
aviation regulator \citep[\FAA{}, \EASA{}, \ANAC{} and other authorities
under \ICAO{} Annex 19;][]{FAA2024,EASA2025} and the industry quality
standard for aerospace, \AS{} \citep{Oschman2017}. Yet, in many
organisations training remains evaluated only at the level of attendance,
written tests or post-class satisfaction surveys --- a posture that limits
training to a compliance obligation and ignores its strategic potential to
reduce waste and reinforce safety culture
\citep{ORconnor2024,Nakamura2024}.

The Kirkpatrick four-level model
\citep{Kirkpatrick2010,Patel2023,Fernandez2024} provides a well-established
framework to evaluate training effectiveness across Reaction, Learning,
Behaviour and Results dimensions. Although broadly cited, its
operationalisation in regulated and high-reliability environments such as
aviation \MRO{} remains under-explored. Recent reviews suggest that fewer than
20\% of published applications go beyond Level~2 (Learning) in any sustained
way, and applications that explicitly connect Level~4 (Results) with
sustainability or cleaner-production outcomes are virtually absent
\citep{ORreilly2024,Peterson2025}.

\subsection{Research questions}
\label{subsec:rqs}

To address this gap, this study investigates the following research
questions (RQs):

\begin{description}
  \item[\textbf{RQ1.}] How can the Kirkpatrick four-level model be
        operationalised within the training process of an \AS{}-certified
        aircraft-engine \MRO{} facility so that it evaluates not only safety
        and regulatory compliance but also resource efficiency and waste
        reduction?
  \item[\textbf{RQ2.}] What measurable changes in safety performance
        (training failures, regulatory non-conformances) and sustainability
        indicators (quality costs, customer satisfaction, turnaround time)
        result from such an implementation over a twelve-month horizon?
  \item[\textbf{RQ3.}] Which managerial and policy interventions can be
        derived from the case so that the framework becomes replicable across
        \MRO{} organisations and useful as guidance for civil aviation
        authorities and industry bodies?
\end{description}

\subsection{Contribution and structure}
\label{subsec:contribution}

To the best of the authors' knowledge, the contribution of this paper is
threefold: \emph{(a)}~it jointly maps Kirkpatrick's four levels to specific
\AS{} clauses, \SMS{} pillars and \ESG{} indicators in an aerospace
maintenance setting, \emph{and (b)}~it reports twelve-month longitudinal
operational outcomes, including a monetised cost-of-quality estimate, from a
single \AS{}-certified engine \MRO{} facility. While prior work has applied
Kirkpatrick to aviation training in isolation and has discussed sustainability
in \MRO{} separately, the joint mapping with longitudinal monetised outcomes
in this regulated setting is, to our knowledge, not yet documented in the
peer-reviewed literature. The third contribution is normative: managerial
implications and policy recommendations directly addressed to \MRO{} managers,
airline quality departments and aviation regulators are derived from the
case. The conceptual mapping is summarised in Fig.~\ref{fig:conceptual}.

\begin{figure*}[!tbp]
  \centering
  \resizebox{0.95\textwidth}{!}{% =============================================================================
% Figure 0 -- Conceptual model (Reviewer comment #8)
% Kirkpatrick levels mapped to AS9100 clauses, SMS pillars, and Sustainability
% =============================================================================
\begin{tikzpicture}[
  every node/.style={font=\small, align=center, inner sep=5pt},
  levelbox/.style={
    rectangle, rounded corners=3pt,
    draw=kpBlue, fill=kpBlueL, thick,
    minimum width=2.6cm, minimum height=1.2cm
  },
  domainbox/.style={
    rectangle, rounded corners=3pt,
    draw=kpGreen, fill=kpGreenL, thick,
    minimum width=3.8cm, minimum height=1.2cm
  },
  sustbox/.style={
    rectangle, rounded corners=3pt,
    draw=kpOrange, fill=kpOrangeL, thick,
    minimum width=3.8cm, minimum height=1.2cm
  },
  arrow/.style={-{Stealth[length=2.5mm]}, thick, gray},
  header/.style={font=\bfseries\small}
]
 % Row positions (1.7cm vertical spacing between rows for breathing room)
 % Header row at y = 6.4
 % L1 at y = 5.1; L2 at y = 3.4; L3 at y = 1.7; L4 at y = 0.0

 % Headers
 \node[header] at (0.0,6.4) {Kirkpatrick};
 \node[header] at (4.2,6.4) {AS9100 / SMS};
 \node[header] at (9.0,6.4) {Sustainability};

 % Level 1
 \node[levelbox] (l1) at (0.0,5.1) {\textbf{L1}\\Reaction};
 \node[domainbox] (a1) at (4.2,5.1) {7.3 Awareness\\7.4 Communication};
 \node[sustbox] (s1) at (9.0,5.1) {Engagement\\(human capital)};
 \draw[arrow] (l1) -- (a1);
 \draw[arrow] (a1) -- (s1);

 % Level 2
 \node[levelbox] (l2) at (0.0,3.4) {\textbf{L2}\\Learning};
 \node[domainbox] (a2) at (4.2,3.4) {7.2 Competence\\(verified knowledge)};
 \node[sustbox] (s2) at (9.0,3.4) {Avoided internal\\failure (scrap)};
 \draw[arrow] (l2) -- (a2);
 \draw[arrow] (a2) -- (s2);

 % Level 3
 \node[levelbox] (l3) at (0.0,1.7) {\textbf{L3}\\Behaviour};
 \node[domainbox] (a3) at (4.2,1.7) {SMS Safety\\Assurance};
 \node[sustbox] (s3) at (9.0,1.7) {Avoided rework\\(energy, labour)};
 \draw[arrow] (l3) -- (a3);
 \draw[arrow] (a3) -- (s3);

 % Level 4
 \node[levelbox] (l4) at (0.0,0.0) {\textbf{L4}\\Results};
 \node[domainbox] (a4) at (4.2,0.0) {Clause 9\\Performance evaluation};
 \node[sustbox] (s4) at (9.0,0.0) {Avoided external failure\\(warranty, CO\textsubscript{2})};
 \draw[arrow] (l4) -- (a4);
 \draw[arrow] (a4) -- (s4);

 % Vertical arrows (Kirkpatrick progression)
 \draw[arrow, dashed] (l1) -- (l2);
 \draw[arrow, dashed] (l2) -- (l3);
 \draw[arrow, dashed] (l3) -- (l4);
\end{tikzpicture}
}
  \caption{Conceptual model linking Kirkpatrick's four levels to AS9100
           clauses, SMS pillars and sustainability outcomes. Vertical dashed
           arrows represent the progression of the Kirkpatrick model;
           horizontal arrows represent the mapping proposed in this study.}
  \label{fig:conceptual}
\end{figure*}

The remainder of the paper is organised as follows.
Section~\ref{sec:literature} reviews the relevant literature on Kirkpatrick,
aviation maintenance training and the safety-sustainability nexus.
Section~\ref{sec:methodology} describes the case-study design and data
collection. Section~\ref{sec:results} presents the results, including the
regulatory mapping, the redesigned training process and the operational
outcomes. Section~\ref{sec:discussion} discusses the findings against the
literature. Section~\ref{sec:managerial} provides managerial implications,
including an economic analysis. Section~\ref{sec:policy} formulates policy
recommendations. Section~\ref{sec:conclusion} concludes, and
Section~\ref{subsec:limitations} synthesises limitations and future research.
